
\medskip


On considère les suites $(v_n)$ et $(w_n)$ définies pour tout entier naturel $n$
par :
\[ \left\{\begin{array}{l c l}
v_0 &=& \ln (4)\\
v_{n+1} &=& \ln \left(- 1 +2\e^{v_n}\right)
\end{array}\right. \qquad \text{et} \qquad w_n = - 1 + \e^{v_n}.\]

On admet que la suite $(v_n)$ est bien définie et strictement positive.

%\medskip
%
%\begin{minipage}{0.6\linewidth}
%\begin{enumerate}
%\item Donner les valeurs exactes de $v_1$ et $w_0$.
%\item 
%	\begin{enumerate}
%		\item Une partie d'une feuille de calcul où figurent les indices et
%les termes des suites $(v_n)$ et $(w_n)$ est reproduite ci-contre.
%		
%Parmi les trois formules ci-dessous, choisir la formule qui,
%saisie dans la cellule B3 puis recopiée vers le bas, permettra
%d'obtenir les valeurs de la suite $(v_n)$ dans la colonne B.
%
%\begin{tabular}{|l |l|}\hline
%Formule 1 &LN($-1 + 2$ * EXP(B2))\\ \hline
%Formule 2 &= LN($-1 + 2$ * EXP(B2))\\ \hline
%Formule 3 &= LN($-1 + 2$ * EXP(A2))\\ \hline
%\end{tabular}
%		\item Conjecturer le sens de variation de la suite $(v_n)$.
%		\item À l'aide d'un raisonnement par récurrence, valider votre
%conjecture concernant le sens de variation de la suite $(v_n)$.
%	\end{enumerate}
%\end{enumerate}
%\end{minipage} \hfill 
%\begin{minipage}{0.36\linewidth}
%{\small\begin{tabular}{c |c|c|c}\hline
%	&A &B &C\\ \hline
%1 	&$n$&$v_n$& $w_n$\\ \hline
%2&0 & 1,38629436 &3\\ \hline
%3& 1& 1,94591015 &6\\ \hline
%4& 2& 2,56494936 &12\\ \hline
%5& 3& 3,21887582 &24\\ \hline
%6& 4& 3,8918203 &48\\ \hline
%7& 5& 4,57471098 &96\\ \hline
%8& 6& 5,26269019 &192\\ \hline
%9& 7& 5,95324333 &384\\ \hline
%10& 8& 6,6450909 7 &768\\ \hline
%11& 9& 7,33758774 &1536\\ \hline
%12& 10& 8,03040956 &3072\\ \hline
%13& 11& 8,72339402 &6144\\ \hline
%14& 12& 9,41645983 &12288\\ \hline
%15& 13& 10,1095663 &24576\\ \hline
%16& 14& 10,8026932 &49152\\ \hline
%17& 15& 11,4958302&98304\\ \hline
%18& 16& 12,1889723&196608\\ \hline
%19& 17&12,8821169& 393216\\ \hline
%\end{tabular}}
%\end{minipage}
%
%\begin{enumerate}[start=3]
%\item
%	\begin{enumerate}
%		\item Démontrer que la suite $(w_n)$ est géométrique.
%		\item En déduire que pour tout entier naturel $n, \: v_n = \ln \left(1 + 3 \times 2^n\right)$.
%		\item Déterminer la limite de la suite $(v_n)$.
%	\end{enumerate}
%\item Justifier que l'algorithme suivant écrit en langage Python renvoie un résultat quel que soit le choix de la valeur du nombre S.
%
%\begin{center}
%\begin{ttfamily}
%\begin{tabular}{|l|}\hline
%from math import*\\
%def seuil(S):\\
%\qquad V=ln(4)\\
%\qquad n=0\\
%\qquad while V < S :\\
%\qquad  \quad n=n+1\\
%\qquad  \quad V=ln(2*exp(V)-1)\\
%\qquad return(n)\\ \hline
%\end{tabular}
%\end{ttfamily}
%\end{center}
%\end{enumerate}

\begin{enumerate}
\item $\bullet~~$$v_1 = \ln \left(- 1 + 2\e^{v_0}\right) = \ln \left(- 1 + 2\e^{\ln 4}\right) = \ln (- 1 + 2 \times 4) = \ln (- 1 + 8)  = \ln 7$ (valeur approchée ligne 4 colonne 3)

$\bullet~~$$w_0 = - 1 + \e^{v_0} = - 1 + \e^{\ln 4} = - 1 + 4 = 3$ (valeur ligne 3 colonne 4)
\item
	\begin{enumerate}
		\item Il faut saisir la formule 2.
		\item On peut penser que la suite $(v_n)$ est croissante
		\item Démonstration par récurrence :

\emph{Initialisation} : on a $v_0 = \ln 4$ et $v_1 = \ln 7$ : la croissance de la fonction ln assure que $v_0 < v_1$.

\emph{Hérédité} : on suppose que pour $n \in \N, \:v_n < v_{n+1}$, alors :

$\e^{v_n} < \e^{v_{n+1}}$ par croissance de la fonction exponentielle

$2\e^{v_n} < 2\e^{v_{n+1}}$ la multiplication par $+ 2$ respecte l'ordre

$- 1 + 2\e^{v_n} < - 1 + 2\e^{v_{n+1}}$ l'addition respecte l'ordre

$\ln \left(- 1 + 2\e^{v_n}\right) < \ln \left(- 1 + 2\e^{v_{n+1}}\right)$ par croissance de la fonction logarithme népérien (on a supposé que la suite $(v_n)$ est bien définie donc que tous les nombres de la forme $- 1 + 2\e^{v_n}$ sont supérieurs à zéro) ; soit finalement 

$v_{n+1} < v_{n+2}$ : l'inégalité est vraie au rang $n + 1$.

Conclusion : l'inégalité est vraie au rang 0 et si elle est vraie au rang $n \in \N$ elle l'est aussi au rang suivant : d'après le principe de récurrence :
\begin{center} quel que soit $n \in \N, \quad v_n < v_{n+1}$ \end{center}

La suite $(v_n)$ est croissante.
	\end{enumerate}
\item
	\begin{enumerate}
		\item Quel que soit $n \in \N$, \: $w_n = - 1 + \e^{v_n}$, donc 
		
$w_{n+1} = - 1 + \e^{v_{n+1}} = - 1 + \e^{\ln \left(- 1 + 2\e^{v_n}\right)} = - 1 - 1 + 2\e^{v_n} = -2 + 2\e^{v_n} = 2\left(- 1 + \e^{v_n}\right)$ : finalement 

Quel que soit $n \in \N$, \quad $w_{n+1} = 2w_n$ : cette égalité montre que la suite $w_n$ est géométrique de raison 2 avec pour premier terme $w_0 = 3$.
		\item On sait qu'alors le terme général $w_n$ est égal à :
		
		\begin{center} $w_n = 3 \times 2^n$ \end{center}
		
Or par définition $w_n = - 1 + \e^{v_n} = 3 \times 2^n \iff \e^{v_n} = 1 + 3 \times 2^n$.

Par croissance de la fonction logarithme népérien : $v_n = \ln \left( 1 + 3 \times 2^n\right)$ pour tout entier naturel.
\item On sait que $\displaystyle\lim_{n \to + \infty} 2^n = + \infty$, donc $\displaystyle\lim_{n \to + \infty} 3 \times 2^n = + \infty$, puis $\displaystyle\lim_{n \to + \infty} 1 + 3 \times 2^n = + \infty$ et enfin $\displaystyle\lim_{n \to + \infty} \ln \left(1 + 3 \times 2^n\right)= + \infty$.

Conclusion : $\displaystyle\lim_{n \to + \infty} v_n = + \infty$.
	\end{enumerate}
\item L'algorithme permet de calculer les termes de la suite $v_n$ à partir de $v_0$. On vient de démontrer que la suite $(v_n)$ n'est pas majorée. Donc quel que soit le choix du nombre $S$, il existe un rang $p$, tel que $v_p > S$ et l'algorithme donnera ce rang $p$.
\end{enumerate}

\medskip

